\section{Findings: Five Speculative Designs for Imagined Future AI Scientific Assistants}
\label{sec:spec-design}


\begin{figure*}
    \centering
    \includegraphics[width=1.0\linewidth]{figures/newspecdesign.png}
    \caption{\textnormal{Five archetypes that are representative of speculative designs that our 12 participants envisioned for how future AI could help lab and field scientists: \textbf{(A)} as the lab's collective knowledge keeper, \textbf{(B)-(C)} distributed lab-wide task status monitor, \textbf{(D)} real-time monitor for scientists' cognitive and physical health, \textbf{(E)} mobile scout for fieldwork, \textbf{(F)} collaborator for hands-on physical chores. \textit{(Image credit: all illustrations were drawn by human artist Lauren Cheng without AI assistance.)}}}
    \Description{A larger illustration shows an overhead view of a lab space with a distributed AI system across multiple rooms labeled (A), (B), and (C), while three illustrations underneath labeled (D), (E), and (F) demonstrate scientists using AI to assist them in their work. In the chemical reagent room (A), a screen located on the lab bench wall houses relevant protocols, notes from colleagues, reminders, images, videos, observations, and warnings specific to the lab. A camera on the device is directed at the work surface. In room (B), cell incubators are highlighted to signify detection of a drop in temperature. In another room (C), a scientist focused on a separate task is alerted about the cell incubators by the screen at her lab bench and by an audio warning. Illustration (D) depicts a scientist wearing a cleanroom suit facing a wall-mounted screen displaying “scanning…,” health information, and suggested tasks. Next to them is a door with a biohazard symbol. In (E), a field scientist working in a remote area deploys stimuli to far away primates via a drone, monitoring their behavior on his device. Illustration (F) shows a woman directing a humanoid robot to pick up a fifty kilogram glass panel used in a muddy terrain tank in her robotics lab.}
    \label{fig:newspecdesign}
\end{figure*}


As each participant was conveying barriers to AI adoption in their workplace, that provided a natural segue into our speculative design activity, where we had them sketch out what an `ideal' future AI assistant would look like for typical tasks. They drew a map of their lab or field space on paper and pointed out how an AI might help. Their design ideas varied depending on the type of task and its location within the lab or out on the field (e.g. lab bench, fume hood, incubation room, surgery room, experiment clean-rooms, outdoors, see Figure ~\ref{fig:science-setting} for examples).

Due to space limitations here, we will not present every single design idea that participants had. Instead we grouped similar-themed ideas together into a set of five archetypal speculative designs and present them here.


\subsection{\textbf{Archetype 1: AI as the lab's collective knowledge keeper}} Across nearly all participants, human limitations in knowledge transfer and documentation emerged as persistent bottlenecks. \rev{Participants} spoke about limitations in their ability to recall and pass down expertise (P1, P2, P3, P4, P5, P6, P7, P9, P10, P12). P3 summed it up as: \textit{``No one likes to do documentation when you're actually working. You just want to work [...] but afterwards you don't remember certain parts. It'd be super helpful if Generative AI could do it for you.''} Experienced scientists struggled to pass down contextual knowledge and expertise to junior colleagues due to fragmented documentation and overwhelming context (P1, P2, P4, P6, P9, P12); and junior scientists struggled to navigate lab-specific, idiosyncratic, tacit practices without clear guidance when a mentor was not nearby (P3, P5, P4).

Some labs maintained shelves of note-filled binders that often went unread, while others had abandoned multiple attempts to create usable computer documentation systems or digital archives. Even when documentation did exist, it was frequently described as dense, confusing, and constantly evolving (P1, P3, P4, P7), exacerbated by each lab and organization having their own custom protocols, tools, equipment: \emph{``There's a lot of tribal knowledge''} (P4).

Thus, in our design sessions participants envisioned context-aware AI systems that could alleviate these gaps by capturing, summarizing, and situating knowledge in real time. For example, P2 imagined a mounted AI camera system that could monitor his robotics experiments through video, digest the streamed footage, and \textit{“provide corrections and live feedback in-context''} for the experiment. He envisioned being able to query the AI, which would then refer back to specific video frames when generating responses.

P1 imagined scenarios where future AI could actively participate in data collection and notetaking: \textit{``There's a ton of digital information that it could automatically take the notes from, `Hey, this thing co-occurred with that thing' [...] or integrating behavioral tracking in my note taking like, `At what point does the rat do this?' Or if I just set up a boundary, `Here's the area I want to keep a track of and here's the rat. Make a note when the rat leaves this area.''}

P12 combined a similar idea with lab-based tacit knowledge-sharing, where context was automatically captured and passed on to colleagues: \textit{``Because we have a couple of people in the lab, if AI could note that I made a mistake there, and for everyone else that has to use this protocol again later, tell them `Please be aware this happened to me. That's how you solve it.' Then, [AI] could pop up and tell you, `Someone had a problem here at this step.'''} See Figure ~\ref{fig:newspecdesign} (A).


\subsection{\textbf{Archetype 2: AI as a distributed lab-wide task status monitor}} Whether in the lab or out in the field, participants frequently identified cognitive limitations that frustrated them, such as difficulty recalling specific procedural knowledge (P1, P2, P3, P4, P5, P10, P11, P12) and lapses in memory or attention in the midst of doing hands-on tasks (P1, P4, P5, P8, P11). These moments of cognitive load were seen as opportunities for AI intervention as a persistent memory aid, context-aware feedback provider, or auxiliary observer.

For example, P1 faced attentional limitations during live experiments with rodents since he needed to juggle multiple simultaneous processes: evaluating data (audio and visual), attending to electrophysiology data, and managing other active tasks, such as \textit{``presenting a certain stimulus at a specific time, taking the rats out, cleaning the arena, setting it back up, putting them back in, baiting a maze [...]''} He also needed to simultaneously attend to equipment failures or apparatus malfunctions, then react quickly. P1 imagined that a camera-based AI assistant would be able to offload some of this burden, acting as an auxiliary monitor---an extra pair of eyes---that can provide attentional scaffolding when cognitive load exceeds human capacity: \textit{``I could totally see it being useful in tracking the things that I might normally be looking at but that I just can't during this time because I'm too focused on this other part or something going wrong.''}

P12 designed for similar attentional limitations when multi-tasking in a cell biology wet-lab with many ongoing tasks and sensitive equipment. Her idea in Figure ~\ref{fig:newspecdesign} (B-C) illustrates experimental protocols distributed across screens in different rooms, and AI monitoring multiple ongoing tasks for her. P12 explained how this system could prevent error by surfacing contextually important information for her at the right time, \textit{``Imagine you walk around and you have a screen there like, `Ah, I'm at this step. I need this and this at this station.' Then, at every screen, every station, you follow the next steps.''} P12 added that this could help prevent careless mistakes; she shared a story where a labmate accidentally left a cell incubator door open, which ended the experiment and delayed their work for weeks. A system could detect the open incubator door in the equipment room (where people usually are not present), illustrated in Figure ~\ref{fig:newspecdesign} (B), and then send an audible alert to a scientist working in a different lab room, as shown in Figure ~\ref{fig:newspecdesign} (C).

P5's hands-free voice-interaction design was informed by his workspace, where access to input devices like computer keyboards was not possible: \textit{``My biggest concern in this space is my hands are often full in some way.}'' During the design session P5 demonstrated how he constructs a rodent brain implant drive under time pressure to check on other tasks:

\begin{quote}
    \textit{``Like, I'm holding this [tetrode wire that is thinner than a human hair], and I finally have the position I want, but I'm thinking, `Hey, is this actually gonna fit?' or I'm not sure how much I have on my alarm left before I need to check on the rats. If I could ask AI on the spot without me having to stop and go to a computer, like a distributed system throughout the lab.''}
\end{quote}

The voice-activated AI monitoring system he envisioned (similar in form to Figure ~\ref{fig:newspecdesign} (C)) aligned with his vision of being able to attend to multiple aspects of the lab concurrently while reducing the friction of task-switching.


\subsection{\textbf{Archetype 3: AI as a real-time monitor for scientists' cognitive and physical health}} Besides monitoring lab tasks, participants also envisioned AI systems that act as status monitors \emph{for the actual scientists} who must do high-stakes and high-intensity work. They brought up concerns that dangerous memory and attention lapses stemmed not only from stressful workloads, but from cognitive and physiological fluctuations in their bodies.

In P4's lab, he and his colleagues were encouraged to self-evaluate one's own headspace and physical condition before certain kinds of intensive work with nuclear materials. Drinking too much caffeine the night before could be reason for pause. P4 added, \textit{``You don't want to be in autopilot [...] I make sure not to do this [task] unless I feel like I'm focused.''} P6, who shares this concern, speculatively designed an AI assistant that could screen for these physical and cognitive fluctuations and alert him accordingly before entering the experiment room: \textit{``If it notices over the past two days, your heart rate was elevated more than normal, it'll warn you.''} See Figure ~\ref{fig:newspecdesign} (D) for an illustration of this idea.

More generally, several \rev{participants} (P1, P4, P5, P10, P12) worried that lapses in memory or careless mistakes could lead to unwanted harm and cascading consequences: \textit{``If your work rhythm gets thrown off, you can forget things [...] if you forget to arm [activate] that thing, you don't collect any data. That's common''} (P5). P5 envisioned querying AI before running an experiment as a quick cognitive status check, \textit{``Hey, I'm about to start, am I forgetting anything?''}

\subsection{\textbf{Archetype 4: AI as a mobile scout for fieldwork}} Participants working in the field envisioned AI assistants that were mobile, lightweight, compatible with remote or rugged environments (P2, P3, P8), and powered by local AI models with offline access (P2, P3). Unlike lab-based researchers, these participants often moved between locations and conducted work that required environmental adaptability.

Because voice interactions with AI could disrupt naturalistic conditions or alter subject behavior (e.g., for animals being observed in close proximity---P8 and P11), they explored alternative modalities in speculative designs. P8, who often collects data at primate sanctuaries, imagined using small autonomous flying drones to reduce the amount of time spent trekking across long distances: \textit{``Oftentimes, there is time wasted to figure out where even the animals are.''} P8 speculated on future capabilities of these drones: \textit{``You can probably efficiently collect a lot of information about where animals are and what they're doing by using a drone.''} Beyond passive observation, P8 envisioned active intervention, such as deploying stimuli or food via drone and recording the resulting behaviors: \textit{``Why couldn't the drone do that for me and just also then film it and see what happens?''} These designs, shown in Figure ~\ref{fig:newspecdesign} (E), illustrate how highly-mobile and non-verbal AI modalities align more with the constraints of field-based research.


\subsection{\textbf{Archetype 5: AI as a collaborator for hands-on physical chores}} Many participants struggled with technical chores that were not core to their research questions but still essential to running their experiments: soldering (P1, P5, P3), epoxying (P8), sawing (P8), wiring (P2, P3, P11), cementing (P1, P5), mixing chemical substances (P1, P6, P9, P12), performing surgical procedures (P4, P6, P10). These skills were often self-taught or learned through trial and error, since they had no formal professional training in those fields. In some cases, recently-acquired skills were only useful for a single experiment. As P8 noted, investing significant time in mastering such skills could slow research momentum, moreso when it can be slow to get feedback on them: \textit{``My PI [principal investigator] would get a chance to look at my [hardware devices] a week after I built them and say, `Okay, this one looks really good. This one isn't sturdy enough. This one isn't wound tight enough' [...] but I'd forgotten what I did already.''}

After surfacing these pain points, participants imagined multimodal AI assistants that could provide scaffolded, live support during such embodied hands-on work. Both P1 and P12 suggested modalities that allowed them to visually learn while they picked up new protocols: \textit{``If I could have some audio thing saying, `What's the next step?' Or have a screen that displays the next step instruction or visual instruction [explaining] what this part is [...]'' (P1).} They shared how audio interaction with AI could provide the best hands-free assistance, alongside visual instructions or diagrams on the screen to serve as persistent references, offering clarity during unfamiliar or error-prone procedures.

Human sensory and physical limitations were also surfaced in the interviews. P12 noted that, \textit{``We have things you cannot see by eye [...] you have to train your muscle memory.''} Multiple participants reported challenges related to the steadiness of their hands (P1, P3, P5, P12) or sustained attention and focus (P1, P4, P12), compounding the difficulty of performing hand-eye coordination tasks that required precision. These participants tended to extend their visions into more embodied, robotic forms of AI assistance; their imagined robots were designed to compensate for human sensory or physical limitations like acting as an extra pair of eyes and hands to help them in these chores.

For example, P10 and P3 envisioned humanoid robots that could help lift, transport, or stabilize materials, whether delicate tiny components like neural cells or heavy equipment like glass panels during experiments (Figure~\ref{fig:newspecdesign} (F)). For P3, this also made certain parts of experiments more accessible: \textit{``An extra pair of hands would help most when you have to take glass panels out [...] it's really heavy.''}


\smallskip \noindent \textbf{\rev{Scientific practitioners} preferred passive AI, with proactive intervention for critical errors.} One recurring theme we observed across design sessions was that most participants (P1, P3, P4, P6, P7, P8, P9, P12) leaned toward passive rather than proactive assistance. Passive assistance stays in the background, offering support only when directly queried by the user. In contrast, proactive assistance engages with users by issuing real-time alerts or suggestions based on human behavior. Some participants were more strongly averse to proactive assistance than others. P9 felt it would be disruptive to her workflow and wanted AI systems to notify users through subtle cues: \textit{``I would prefer the robot to raise its hand or flash until I touch it. I don't want it to interact and directly talk to me like, `Hey, you missed that!' That would drive me nuts.''}

Modalities of AI intervention that came up most frequently in sessions were wall-mounted screens and hands-free, voice-interactive systems. Notably, wall-mounted screens \emph{without audio or voice controls} were a preferred modality for participants who were concerned that audio could be a potential confound or distraction when running experiments (P1, P11, P12). In these scenarios, participants preferred AI assistants to provide suggestions, instructions, and warnings via text on large mounted screens, or intervene through minimal visual cues, like flashing red lights (P9). More generally, P11 did not want future AI systems to deprive her of the opportunity to \textit{``do the actual work [myself]}.''

That said, some (P1, P2, P5, P10, P11, P12) favored a mix of passive and proactive. They emphasized that although passive systems were preferred, they wanted AI to intervene when a mistake was imminent, especially in tasks where mistakes could be costly or hard to detect (P2, P12). P5 wanted control over this threshold of AI intervention: \textit{``Ideally, there'd be different preferences. For one, because you obviously don't want [AI] to every 30 seconds go, `Hey, do this. Hey, do that.' But in emergency situations, like, `Hey, you forgot to arm a system,' I want that.''}